\documentclass[10pt,a4paper]{amsart}
\usepackage[margin=19mm]{geometry}
\usepackage{amsmath,amssymb,mathtools}
\usepackage[scr=rsfs,cal=euler]{mathalfa}
\usepackage{xfrac}
\usepackage[unicode,hidelinks,bookmarks=false]{hyperref}
\newcommand{\ie}{i.e.\ }
\newcommand{\cf}{cf.\ }
\newcommand{\resp}{respectively\ }
\newcommand{\Subsection}{\subsection}
\providecommand{\nobreakdash}{\nobreak\hbox{-}}
\setlength{\emergencystretch}{2em}
\multlinegap0pt
\usepackage{bm}
\usepackage{bbm}
\usepackage{dsfont}
\usepackage{xspace}
\usepackage{cancel}
\usepackage{mathtools}
\usepackage[shortlabels]{enumitem}
\usepackage{amsthm}
\makeatletter
\@ifundefined{thm}{\newtheorem{thm}{Thm}}{}
\@ifundefined{lem}{\newtheorem{lem}[thm]{Lemma}}{}
\makeatother
\newcommand{\e}{\epsilon}
\newcommand{\ra}{\rightarrow}
\newcommand{\s}{\sigma}
\renewcommand{\k}{\kappa}
\title[Blaschke-Santal\'o type inequalities and quermassintegral in]{Blaschke-Santal\'o type inequalities and quermassintegral inequalities in space forms}
\author{Yingxiang Hu, Haizhong Li}
\date{Source: arXiv:2207.04281v1 (CC BY 4.0)}
\begin{document}
\maketitle

\noindent\emph{Source and licence.} Blaschke-Santal\'o type inequalities and quermassintegral inequalities in space forms. Version arXiv:2207.04281v1 (CC BY 4.0), distributed under the Creative Commons Attribution 4.0 International licence, as recorded in the arXiv licence metadata for this exact version and shown on the abstract page https://arxiv.org/abs/2207.04281v1. Full rights notice: licenses/arxiv-2207.04281-CC-BY-4.0.txt.\par\smallskip

\noindent\textit{Editorial scope.} Counted unit: the Lem whose source label is \texttt{s4:lem-identity} in the article (source0.tex lines 670--681, source label \texttt{s4:lem-identity}), delivered together with the article's own complete original proof of it (source0.tex lines 682--709, one \texttt{proof} environment of 28 lines). No mathematics is written, repaired or re-derived: statement and proof are the article's own text line for line. Required context retained uncounted: s2:duality-sphere (lines 446--448); s2:duality-hyperbolic (lines 475--477); s2:duality-deSitter (lines 509--511); s2:curvature-flow (lines 523--529); s2:dual-function (lines 546--548); s2:dual-flow (lines 550--556); s3:variational-formula (lines 596--606). Every definition, display and label of those lines is retained unchanged. Omitted from this package and not redistributed: 1--445, 449--474, 478--508, 512--522, 530--545, 549--549, 557--595, 607--669, 710--1638. Nothing inside a retained excerpt is omitted or altered: every retained body is delivered exactly as the article's lines are written, so no omission receipt of any kind accompanies this package. Citations: the retained excerpts carry no citation command at all, so no bibliography entry of the article is transcribed and no reference is redistributed. Reference adaptation: none was needed and none was made; every reference inside the delivered unit resolves inside it, so no unresolved reference or citation is produced. Printed slips kept literal: no printed word, symbol, hypothesis or number of the counted statement or of its proof was altered. Layout and class: the article's own class is \texttt{amsart} with packages including enumerate, amsrefs, amsfonts, amsmath, amssymb, amscd, amsthm, bm, cancel, mathrsfs, url, graphicx, hyperref, multicol; the wrapper is the lane's pinned \texttt{amsart} editorial wrapper at 10pt on A4, which restates the theorem-like environments the retained text needs and adds the article's own macro definitions that the retained text needs (\texttt{\textbackslash e}, \texttt{\textbackslash k}, \texttt{\textbackslash ra}, \texttt{\textbackslash s}), with extra wrapper packages (bm, bbm, dsfont, xspace, cancel, mathtools, enumitem). The delivered numbering is the unit's own. Assets: no retained span contains a figure environment, a graphics-inclusion command, a PDF-page inclusion, a table or a TikZ picture; no image file is copied into this package, the package declares no asset dependency and no image file is redistributed with it. Licence: version arXiv:2207.04281v1 of this article is distributed under the Creative Commons Attribution 4.0 International licence, as recorded in the arXiv licence metadata for this exact version and shown on the abstract page https://arxiv.org/abs/2207.04281v1. A full rights notice is delivered at licenses/arxiv-2207.04281-CC-BY-4.0.txt. Characters: every retained excerpt is pure ASCII, so the delivered engine, XeLaTeX with its default Latin Modern text font, reports no missing character.
	\begin{align}\label{s2:duality-sphere}
	h_{ij}=h^\ast_{ij}=\langle X^\ast_i,X_j\rangle, \quad \k_i^\ast=\k_i^{-1},
	\end{align}

	\begin{align}\label{s2:duality-hyperbolic}
	h_{ij}=h^\ast_{ij}=\langle X^\ast_i,X_j\rangle, \quad \k_i^\ast=\k_i^{-1}.
	\end{align}

	\begin{align}\label{s2:duality-deSitter}
	h_{ij} =h^\ast_{ij}=\langle X_i,X^\ast_j\rangle, \quad 	\k^\ast_i=\k_i^{-1},
	\end{align}	

\begin{align}\label{s2:curvature-flow}
\left\{
\begin{aligned}
\frac{\partial}{\partial t}X=&-\s \mathcal{F}\nu, \\
X(\cdot,0)=&X_0(\cdot).
\end{aligned}\right.
\end{align}

\begin{align}\label{s2:dual-function}
\mathcal{F}^\ast(\mathcal{W}^\ast):=-\mathcal{F}(\mathcal{W}).
\end{align}

\begin{align}\label{s2:dual-flow}
\left\{
\begin{aligned}
\frac{\partial}{\partial t}X^\ast=&-\s^{\ast} \mathcal{F}^\ast \nu^\ast-(b^\ast)^{kl}\mathcal{F}^\ast_{k}X^\ast_l, \\
X^\ast(\cdot,0)=&X^\ast_0(\cdot).
\end{aligned}\right.
\end{align}

\begin{lem}
	Assume that one of the following holds:
	\begin{enumerate}[(i)]
		\item Let $K_t$ be a family of smooth bounded domains in $\mathbb N^{n+1}(\e)$ such that $\partial K_t$ is a family of smooth closed hypersurfaces. 
		\item Let $K_t$ be a family of smooth bounded domains in $\mathbb S^{n,1}_{+}$ such that $\partial K_t$ is a family of smooth closed spacelike hypersurfaces.
    \end{enumerate}
    Let $\frac{\partial}{\partial t}X$ be the variational vector field of $\partial K_t$. Then 
	\begin{align}\label{s3:variational-formula}
	\frac{d}{dt}W_k(K_t)=\frac{n+1-k}{n+1}\int_{\partial K_t} \langle \frac{\partial}{\partial t}X,\nu \rangle E_k d\mu_t, \quad k=0,1,\cdots,n.
	\end{align}	
\end{lem}

\begin{lem}\label{s4:lem-identity}
	Assume that one of the following holds:
	\begin{enumerate}[(i)]
		\item Let $K_t$ be a family of smooth bounded domain with strictly convex boundary $\partial K_t$ in $\mathbb S^{n+1}$, and $K_t^\ast$ the dual body of $K_t$ in $\mathbb S^{n+1}$;
		\item Let $K_t$ be a family of smooth bounded domain with strictly convex boundary $\partial K_t$ in $\mathbb H^{n+1}$, and $K_t^\ast$ the dual body of $K_t$ in $\mathbb S^{n,1}_{+}$;
		\item Let $K_t$ be a family of smooth bounded domain with strictly convex and spacelike boundary $\partial K_t$ in $\mathbb S^{n,1}_{+}$, and $K_t^\ast$ the dual body of $K_t$ in $\mathbb H^{n+1}$.
\end{enumerate}	
	    Assume that $X:M_0 \times [0,T) \ra N$ is a family of smooth embeddings satisfying \eqref{s2:curvature-flow} such that $X_t(M_0)=\partial K_t$, then there holds
	\begin{align}\label{s4:evol-identity}
	\frac{d}{dt}\left[ (k+1)W_k(K_t)+(n+1-k)W_{n-k}(K_t^\ast)\right]=0, \quad k=0,\cdots,n.
	\end{align}
\end{lem}

\begin{proof}
	Along the flow \eqref{s2:curvature-flow} in $N$, i.e.,
	\begin{align*}
	\frac{\partial}{\partial t}X=-\s \mathcal{F}\nu,
	\end{align*}
	it follows from \eqref{s3:variational-formula} that
	\begin{align}\label{s4:quermassintegral-direct-flow}
	\frac{d}{dt}W_k(K_t)=&\frac{n+1-k}{n+1}\int_{\partial K_t} E_k(\k) \langle \frac{\partial}{\partial t}X,\nu\rangle d\mu_t\nonumber \\ 
	=&-\frac{n+1-k}{n+1}\int_{M_t} E_k(\k) \mathcal{F} d\mu_t,
	\end{align}
	where $\langle \cdot,\cdot\rangle$ denotes the Euclidean inner product for $\mathbb S^{n+1}$ and the Minkowski inner product for $\mathbb H^{n+1}$ and $\mathbb S^{n,1}_{+}$, respectively.
	 
	Corresponding to the flow \eqref{s2:curvature-flow}, the dual flow \eqref{s2:dual-flow} in the dual ambient space $N^\ast$ is 
	\begin{align}
	\frac{\partial}{\partial t}X^\ast=-\s^\ast \mathcal{F}^\ast \nu^\ast-(b^\ast)^{kl}\mathcal{F}^\ast_k X^\ast_l.
	\end{align}
	It also follows from \eqref{s3:variational-formula} that
	\begin{align}\label{s4:quermassintegral-dual-flow}
	\frac{d}{dt}W_{n-k}(K^\ast_t)=&\frac{n+1-(n-k)}{n+1}\int_{\partial K^\ast_t} E_{n-k}(\k^\ast) \langle \frac{\partial}{\partial t}X^\ast, \nu^\ast \rangle d\mu^\ast_t \nonumber\\
	=&-\frac{n+1-(n-k)}{n+1}\int_{\partial K^\ast_t} E_{n-k}(\k^\ast) \mathcal{F}^\ast d\mu^\ast_t \nonumber\\
	=&\frac{k+1}{n+1}\int_{\partial K_t} E_k(\k) \mathcal{F} d\mu_t,
	\end{align}
	where we used \eqref{s2:dual-function} and 
	\begin{align*}
	E_{n-k}(\k^\ast)=\frac{E_k(\k)}{E_{n}(\k)}, \quad d\mu^\ast_t=E_n(\k) d\mu_t
	\end{align*}
	due to the duality relations \eqref{s2:duality-sphere}, \eqref{s2:duality-hyperbolic} and \eqref{s2:duality-deSitter}. Therefore, the identity \eqref{s4:evol-identity} follows from \eqref{s4:quermassintegral-direct-flow} and \eqref{s4:quermassintegral-dual-flow}.	
\end{proof}

\end{document}
